\section{矩阵微积分：手动推导梯度}

\subsection{从单变量到多变量}

Manning 教授给出了一个关键的思维口诀：

\begin{importantbox}{矩阵微积分的核心口诀}
\textbf{``Multivariable calculus is just like single-variable calculus if you use matrices.''}（多变量微积分就是单变量微积分——只不过用矩阵代替标量。）
\end{importantbox}

\begin{figure}[H]
\centering
\includegraphics[width=0.95\textwidth]{slides-images/slide-010.jpg}
\caption{手动计算梯度：矩阵微积分 = 全向量化的梯度计算，比逐元素计算快得多\protect\footnotemark}
\end{figure}
\footnotetext{Slides 第11页。}

回顾单变量微积分的基本例子：$f(x) = x^3$，导数 $f'(x) = 3x^2$。导数衡量的是\textbf{斜率}——输入的微小变化如何被放大到输出上。例如在 $x = 4$ 处，$f'(4) = 48$，意味着 $x$ 从 4 变到 4.01 时，$f(x)$ 大约增加 $48 \times 0.01 = 0.48$（即从 64 变为约 64.48）。

\subsection{梯度与 Jacobian 矩阵}

当函数有 $n$ 个输入和 1 个输出时（$f: \mathbb{R}^n \to \mathbb{R}$），梯度是一个向量：

$$\nabla f = \begin{pmatrix} \frac{\partial f}{\partial x_1} \\ \frac{\partial f}{\partial x_2} \\ \vdots \\ \frac{\partial f}{\partial x_n} \end{pmatrix}$$

当函数有 $n$ 个输入和 $m$ 个输出时（$f: \mathbb{R}^n \to \mathbb{R}^m$），偏导数构成一个 $m \times n$ 的矩阵，称为 \textbf{Jacobian 矩阵}：

$$\frac{\partial \mathbf{f}}{\partial \mathbf{x}} = \begin{pmatrix} \frac{\partial f_1}{\partial x_1} & \cdots & \frac{\partial f_1}{\partial x_n} \\ \vdots & \ddots & \vdots \\ \frac{\partial f_m}{\partial x_1} & \cdots & \frac{\partial f_m}{\partial x_n} \end{pmatrix}$$

\begin{figure}[H]
\centering
\includegraphics[width=0.95\textwidth]{slides-images/slide-014.jpg}
\caption{Jacobian 矩阵的定义：$m$ 个输出对 $n$ 个输入的偏导数排列成 $m \times n$ 矩阵\protect\footnotemark}
\end{figure}
\footnotetext{Slides 第15页。}

\begin{knowledgebox}{Jacobian 在神经网络中的角色}
神经网络的每一层都是一个多输入多输出的函数。例如一个隐藏层 $h = f(Wx + b)$ 将 $m$ 维输入映射到 $n$ 维输出。该层的 Jacobian 就是一个 $n \times m$ 的矩阵，描述了每个输出维度对每个输入维度的敏感度。
\end{knowledgebox}

\subsection{链式法则的矩阵形式}

对于复合函数，链式法则变为 Jacobian 矩阵的\textbf{乘法}。如果 $h = f(g(x))$，则：

$$\frac{\partial h}{\partial x} = \frac{\partial h}{\partial g} \cdot \frac{\partial g}{\partial x}$$

这与单变量的链式法则 $\frac{dz}{dx} = \frac{dz}{dy} \cdot \frac{dy}{dx}$ 形式完全一致——只是标量乘法变成了矩阵乘法。

\begin{figure}[H]
\centering
\includegraphics[width=0.95\textwidth]{slides-images/slide-017.jpg}
\caption{链式法则：复合函数的 Jacobian = 各层 Jacobian 的矩阵乘积\protect\footnotemark}
\end{figure}
\footnotetext{Slides 第18页。}

\subsection{逐元素激活函数的 Jacobian}

对于逐元素激活函数 $h = f(z)$（其中 $h_i = f(z_i)$），其 Jacobian 是一个\textbf{对角矩阵}：

\begin{figure}[H]
\centering
\includegraphics[width=0.95\textwidth]{slides-images/slide-020.jpg}
\caption{逐元素激活函数的 Jacobian：$\frac{\partial h}{\partial z} = \text{diag}(f'(z))$，是对角矩阵\protect\footnotemark}
\end{figure}
\footnotetext{Slides 第21页。}

$$\left(\frac{\partial \mathbf{h}}{\partial \mathbf{z}}\right)_{ij} = \begin{cases} f'(z_i) & \text{if } i = j \\ 0 & \text{otherwise} \end{cases}$$

即：
$$\frac{\partial \mathbf{h}}{\partial \mathbf{z}} = \text{diag}(f'(\mathbf{z}))$$

\begin{importantbox}{对角结构的直觉}
为什么是对角矩阵？因为激活函数是\textbf{逐元素}应用的：$h_i$ 只依赖于 $z_i$，不依赖于 $z_j$（$j \neq i$）。因此非对角元素（$i \neq j$）的偏导数为零。
\end{importantbox}

\subsection{线性层的 Jacobian}

对于线性变换 $z = Wx + b$：

$$\frac{\partial z}{\partial x} = W, \qquad \frac{\partial z}{\partial b} = I$$

类比单变量：$y = wx + b$ 的导数关于 $x$ 是 $w$，关于 $b$ 是 $1$。矩阵版本中，$w$ 变为 $W$，$1$ 变为单位矩阵 $I$。

对于点积 $s = u^T h$：

$$\frac{\partial s}{\partial h} = u^T$$

\begin{figure}[H]
\centering
\includegraphics[width=0.95\textwidth]{slides-images/slide-022.jpg}
\caption{常用 Jacobian 汇总：$\frac{\partial(Wx+b)}{\partial x} = W$，$\frac{\partial(Wx+b)}{\partial b} = I$，$\frac{\partial(u^Th)}{\partial h} = u^T$\protect\footnotemark}
\end{figure}
\footnotetext{Slides 第23页。}

\subsection{NER 网络的完整梯度推导}

现在将所有 Jacobian 组合起来，对 NER 示例网络推导梯度。网络结构为：

$$s = u^T h, \quad h = f(z), \quad z = Wx + b$$

\begin{figure}[H]
\centering
\includegraphics[width=0.95\textwidth]{slides-images/slide-030.jpg}
\caption{应用链式法则：$\frac{\partial s}{\partial b} = \frac{\partial s}{\partial h} \cdot \frac{\partial h}{\partial z} \cdot \frac{\partial z}{\partial b}$\protect\footnotemark}
\end{figure}
\footnotetext{Slides 第31页。}

\textbf{计算 $\frac{\partial s}{\partial b}$}：

$$\frac{\partial s}{\partial b} = \frac{\partial s}{\partial h} \cdot \frac{\partial h}{\partial z} \cdot \frac{\partial z}{\partial b} = u^T \cdot \text{diag}(f'(z)) \cdot I$$

化简后：
$$\frac{\partial s}{\partial b} = u^T \circ f'(z)$$

其中 $\circ$ 表示 \textbf{Hadamard 乘积}（逐元素乘法）。

\begin{knowledgebox}{Hadamard 乘积}
Hadamard 乘积（记作 $\circ$ 或 $\odot$）是将两个相同维度的向量/矩阵\textbf{逐元素相乘}。不同于点积（结果是标量），Hadamard 乘积的结果维度与输入相同。在神经网络的梯度计算中，它出现在对角 Jacobian 乘以向量时——等价于两个向量的逐元素乘法。
\end{knowledgebox}

\textbf{上游梯度 $\delta$ 的概念}：

注意到计算 $\frac{\partial s}{\partial b}$ 和 $\frac{\partial s}{\partial W}$ 时，前两项是共享的：

$$\delta = \frac{\partial s}{\partial h} \cdot \frac{\partial h}{\partial z} = u^T \circ f'(z)$$

$\delta$ 称为\textbf{上游梯度}（upstream gradient）或\textbf{误差信号}（error signal）。它可以计算一次、复用多次，这是反向传播算法效率的关键。

\begin{figure}[H]
\centering
\includegraphics[width=0.95\textwidth]{slides-images/slide-035.jpg}
\caption{$\delta$（上游梯度/误差信号）的复用：$\frac{\partial s}{\partial b} = \delta$，$\frac{\partial s}{\partial W} = \delta^T x^T$\protect\footnotemark}
\end{figure}
\footnotetext{Slides 第36页。}

于是：
\begin{itemize}
    \item $\frac{\partial s}{\partial b} = \delta$（因为 $\frac{\partial z}{\partial b} = I$）
    \item $\frac{\partial s}{\partial W} = \delta^T x^T$（外积形式）
\end{itemize}

\textbf{计算 $\frac{\partial s}{\partial W}$}：

\begin{figure}[H]
\centering
\includegraphics[width=0.95\textwidth]{slides-images/slide-040.jpg}
\caption{对矩阵求导的形状问题：$W \in \mathbb{R}^{n \times m}$，严格的 Jacobian 是 $1 \times nm$ 行向量\protect\footnotemark}
\end{figure}
\footnotetext{Slides 第41页。}

按照严格的 Jacobian 定义，$\frac{\partial s}{\partial W}$ 应该是一个 $1 \times nm$ 的行向量（因为 $s$ 是标量，$W$ 有 $nm$ 个参数）。但在实际工程中，我们采用\textbf{形状约定}（shape convention），将梯度整形为与参数\textbf{相同的形状}，以便直接做 SGD 更新 $W^{\text{new}} = W^{\text{old}} - \alpha \frac{\partial s}{\partial W}$。

最终结果：
$$\frac{\partial s}{\partial W} = \delta^T x^T \in \mathbb{R}^{n \times m}$$

这是 $\delta^T$（$n \times 1$）与 $x^T$（$1 \times m$）的\textbf{外积}。

\begin{warningbox}{Jacobian 形式 vs 形状约定}
数学上，Jacobian 矩阵有严格的定义和形状，链式法则在 Jacobian 形式下是矩阵乘法。但在工程实现中，为了方便 SGD 更新，我们通常将梯度``reshape''为与参数相同的形状。这两种约定混用可能导致困惑——Manning 教授建议：可以先用 Jacobian 做正确的数学推导，最后再 reshape 结果；或者始终按形状约定做，但要灵活使用转置。
\end{warningbox}

\begin{figure}[H]
\centering
\includegraphics[width=0.95\textwidth]{slides-images/slide-044.jpg}
\caption{$\frac{\partial s}{\partial W}$ 的直觉解释：权重 $W_{23}$ 只连接输入 $x_3$ 到 $h_2$，所以其梯度只依赖于 $\delta_2$ 和 $x_3$\protect\footnotemark}
\end{figure}
\footnotetext{Slides 第45页。}

\subsection{外积梯度的直觉}

为什么 $\frac{\partial s}{\partial W} = \delta^T x^T$ 是对的？考虑单个权重 $W_{ij}$：它只参与计算 $z_i = \sum_k W_{ik} x_k + b_i$，即只连接输入 $x_j$ 到隐藏层 $z_i$。因此：

$$\frac{\partial s}{\partial W_{ij}} = \delta_i \cdot x_j$$

这恰好是外积 $\delta^T x^T$ 的第 $(i,j)$ 个元素。

\subsection{本章小结}

矩阵微积分的核心是：用 Jacobian 矩阵表示多输入多输出函数的导数，用矩阵乘法表示链式法则。对于神经网络的常见操作（线性层、激活函数、点积），Jacobian 都有简洁的封闭形式。关键概念是\textbf{上游梯度 $\delta$}——它在不同参数的梯度计算中被复用，避免重复计算。

%% ========== Part 6: 反向传播算法 ==========

